\documentclass[twoside,11pt,openright]{report}

% \usepackage[latin1]{inputenc}
\usepackage[english]{babel}
\usepackage{a4}
\usepackage{latexsym}
\usepackage{amssymb}
\usepackage{amsmath}
\usepackage{epsfig}
\usepackage[T1]{fontenc}
% \usepackage{mathptmx} % Use Times New Romans as font?
\usepackage{color}
\usepackage{epstopdf}
\usepackage{microtype}
\usepackage[hidelinks]{hyperref}
\usepackage[useregional]{datetime2}
\DTMlangsetup[en-US]{showdayofmonth=false}
\usepackage{lipsum}

\renewcommand*\sfdefault{lmss}
\renewcommand*\ttdefault{txtt}

% Custom packages added:

\usepackage{fullpage}
\setlength\marginparwidth{0.7in}

\usepackage{todonotes} % \todo command
\usepackage{proof} % \infer command
\usepackage{tikz-cd}
\usepackage{xfrac} % \sfrac - sideways fraction

% Custom commands

\newcommand*{\figref}[1]{Figure~\ref{#1}}

\newcommand*{\set}[1]{\left\{#1\right\}}
\newcommand*{\term}[1]{\textcolor{blue!80!black}{#1}}
\newcommand*{\type}[1]{\textcolor{magenta!90!black}{#1}}
\newcommand*{\dependenttype}[1]{\textcolor{magenta!70!black}{#1}}
\newcommand*{\container}[1]{\textcolor{red}{#1}}
\newcommand*{\containerpair}[2]{\textcolor{red}{(}#1\,\textcolor{red}{,}\,#2\textcolor{red}{)}}
\newcommand*{\universe}[1]{\textcolor{orange!80!black}{#1}}
\newcommand*{\unit}{\type{\textbf{1}}}
\newcommand*{\empt}{\type{\textbf{0}}}
\newcommand*{\coalg}[2]{#1\texttt{-}#2}

\newcommand*{\function}[1]{\textcolor{blue!60!black}{\ensuremath{\mathtt{#1}}}}
\newcommand*{\typeformer}[1]{\ensuremath{\mathtt{#1}}}
\newcommand*{\functor}[1]{\ensuremath{\mathbf{\mathtt{#1}}}}

\newcommand*{\ssfrac}[2]{ {\LARGE \text{\sfrac{\ensuremath{#1}}{\ensuremath{#2}}}}}

\begin{document}

%%%%%%%%%%%%%%%%%%%%%%%%%%%%%%%%%%%%%%%%%%%%%%%%%%%%%%%%%%%%%%%%%%%%%%%

\pagestyle{empty} 
\pagenumbering{roman} 
\vspace*{\fill}\noindent{\rule{\linewidth}{1mm}\\[4ex]
{\Huge\sf M-types and Coinduction in HoTT and Cubical Type Theory}\\[2ex]
{\huge\sf Lasse Letager Hansen, 201912345}\\[2ex]
\noindent\rule{\linewidth}{1mm}\\[4ex]
\noindent{\Large\sf Master's Thesis, Computer Science\\[1ex] 
\today \\[1ex] Advisor: Bas Spitters\\[15ex]}\\[\fill]}
\epsfig{file=logo.eps}\clearpage

%%%%%%%%%%%%%%%%%%%%%%%%%%%%%%%%%%%%%%%%%%%%%%%%%%%%%%%%%%%%%%%%%%%%%%%

\pagestyle{plain}
\chapter*{Abstract}
\addcontentsline{toc}{chapter}{Abstract}

\todo[inline]{in English\dots}

\chapter*{Resum\'e}
\addcontentsline{toc}{chapter}{Resum\'e}

\todo[inline]{in Danish\dots}

\chapter*{Acknowledgments}
\addcontentsline{toc}{chapter}{Acknowledgments}

\todo[inline]{\dots}

\vspace{2ex}
\begin{flushright}
  \emph{Lasse Letager Hansen,}\\
  \emph{Aarhus, \today.}
\end{flushright}

\tableofcontents
\cleardoublepage
\pagenumbering{arabic}
\setcounter{secnumdepth}{2}

%%%%%%%%%%%%%%%%%%%%%%%%%%%%%%%%%%%%%%%%%%%%%%%%%%%%%%%%%%%%%%%%%%%%%%%

\chapter{Introduction}
\label{ch:intro}

\todo[inline]{motivate and explain the problem to be addressed}

\todo[inline]{example of a citation: \cite{amin2018}}
\todo[inline]{get your bibtex entries from \url{https://dblp.org/}}

%%%%%%%%%%%%%%%%%%%%%%%%%%%%%%%%%%%%%%%%%%%%%%%%%%%%%%%%%%%%%%%%%%%%%%%

\chapter{M-types}
\label{ch:m-types}

\section{Containers / Signatures}
A Container (or Signature) is a pair \(\container{S} = (\type{A},\type{B})\) of types \(\vdash \type{A} : \universe{\mathcal{U}}\) and \(\term{a} : \type{A} \vdash \dependenttype{B}(\term{a}) : \universe{\mathcal{U}}\). From a container we can define a polynomial functor, defined for objects (types) as
\begin{equation}
  \begin{aligned}
    \functor{P}_{\container{S}} &: \universe{\mathcal{U}} \rightarrow \universe{\mathcal{U}} \\
    \functor{P}(\type{X}) := \functor{P}_{\container{S}}(\type{X}) &= \sum_{\term{a} : \type{A}} \dependenttype{B}(\term{a}) \rightarrow \type{X}
  \end{aligned}
\end{equation}
and for a function \(\term{f} : \type{X} \rightarrow \type{Y}\) as
\begin{equation}
  \begin{aligned}
    \functor{P}\term{f} &: \functor{P}\type{X} \rightarrow \functor{P}\type{Y} \\
    \functor{P}\term{f}(\term{a},\term{g}) &= (a, \term{f} \circ \term{g})
  \end{aligned}
\end{equation}
As an example lets look at type for streams over the type \(\type{A}\), defined by the container \(\container{S} = (\type{A},\lambda\,\_, \unit)\), applying the polynomial functor we get
\begin{equation}
  \functor{P}_{\container{S}}(\type{X}) = \sum_{\term{a} : \type{A}} \unit \rightarrow \type{X}
\end{equation}
since we are working in a Category with exponentials we get \(\unit \rightarrow \type{X} \equiv \type{X}^{\unit} \equiv \type{X}\), furthermore \(\unit\) and \(\type{X}\) does not depend on \(\type{A}\) here, so this will be equivalent to the definition
\begin{equation}
  \functor{P}_{\container{S}}(\type{X}) = \type{A} \times \type{X}
\end{equation}
Now we define the coalgebra for this functor with type
\begin{equation}
  \typeformer{Coalg}_{\container{S}} = \sum_{\type{C} : \universe{\mathcal{U}}} \type{C} \rightarrow \functor{P}\type{C}
\end{equation}
we denote a coalgera of \(\type{C}\) and \(\term{\gamma}\) as \(\coalg{\type{C}}{\term{\gamma}}\). The coalgebra morphisms are defined as
\begin{equation}
  \begin{aligned}
    \_\function{\Rightarrow}\_ &: \typeformer{Coalg}_{\container{S}} \rightarrow \typeformer{Coalg}_{\container{S}} \\
    \coalg{\type{C}}{\term{\gamma}}~\function{\Rightarrow}~\coalg{\type{D}}{\term{\delta}} &= \sum_{\term{f} : \type{C} \rightarrow \type{D}} \term{\delta} \circ \term{f} = \functor{P} \term{f} \circ \term{\gamma}
  \end{aligned}
\end{equation}
\(\mathtt{M}\)-types can now be defined from a container \(\container{S}\) as the type \(\typeformer{M}_{\container{S}}\) such that the coalgebra for \(\typeformer{M}_{\container{S}}\) and \(\function{out} : \typeformer{M}_{\container{S}} \rightarrow \functor{P}_{\container{S}}(\typeformer{M}_{\container{S}})\) fulfills the property
\begin{equation}
  \typeformer{Final}_{\container{S}} := \sum_{(\coalg{\type{X}}{\term{\rho}} : \mathtt{Coalg}_{\container{S}})} \prod_{(\coalg{\type{C}}{\term{\term{\gamma}}} : \mathtt{Coalg}_{\container{S}})} \mathtt{isContr}(\coalg{\type{C}}{\term{\gamma}} \Rightarrow \coalg{\type{X}}{\term{\rho}})
\end{equation}
that is \(\prod_{(\coalg{\type{C}}{\term{\gamma}} : \mathtt{Coalg}_{\container{S}})} \mathtt{isContr}(\coalg{\type{C}}{\term{\gamma}} \Rightarrow \coalg{\typeformer{M}_{\container{S}}}{\function{out}})\). We denote this construction of the M-type as \(\typeformer{M}_{\containerpair{\type{A}}{\type{B}}}\) or \(\typeformer{M}_{\container{S}}\) or just \(\typeformer{M}\) when the Container is clear from the context.
\\ \\
If we continue our example for streams this will give us the \(\mathtt{M}\)-type, we can see that \(\functor{P}_{\container{S}}(\typeformer{M}) = \type{A} \times \typeformer{M}\), meaning we have the following diagram,
\begin{figure}[h]
  \centering
  \begin{tikzcd}[row sep=large,column sep=large]
    \type{A} & \type{A} \times \typeformer{M}_{\containerpair{A}{\lambda\,\_,\unit}} \ar[l,"\pi_1"'] \ar[r,"\pi_2"] & \typeformer{M}_{\containerpair{A}{\lambda\,\_,\unit}} \\
    & \typeformer{M}_{\containerpair{A}{\lambda\,\_,\unit}} \ar[ul,"\function{hd}"] \ar[u,"\function{out}"'] \ar[ur,"\function{tl}"']
  \end{tikzcd}
  \caption{M-types of streams}
\end{figure}
where \(\function{out}\) is an isomorphism (because of the finality of the coalgebra), with inverse \(\function{in} : \functor{P}_{\container{S}}(\typeformer{M}) \rightarrow \typeformer{M}\). We now have a semantic for the rules we would expect for streams, if we let \(\function{cons} = \function{in}\) and \(\typeformer{stream}~\type{A} = \typeformer{M}_{\containerpair{\type{A}}{\lambda\,\_,\unit}}\),
\begin{equation}
  \infer[\mathtt{E}_{\function{hd}}]{\function{hd}~\term{s} : \type{A}}{\type{A} : \universe{\mathcal{U}} & \term{s} : \typeformer{stream}~\type{A}}
\end{equation}
\begin{equation}
  \infer[\mathtt{E}_{\function{tl}}]{\function{tl}~\term{s} : \typeformer{stream}~\type{A}}{\type{A} : \universe{\mathcal{U}} & \term{s} : \typeformer{stream}~\type{A}}
\end{equation}
\begin{equation}
  \infer[\mathtt{I}_{\function{cons}}]{\term{\function{cons}}~\term{x}~\term{xs} : \typeformer{stream}~\type{A}}{\type{A} : \universe{\mathcal{U}} & \term{x} : \type{A} & \term{xs} : \typeformer{stream}~\type{A}}
\end{equation}

\section{ITrees as \texttt{M}-types}
We want the following rules for ITrees
\begin{equation}
  \infer[\mathtt{I}_{\function{Ret}}]{\function{Ret}~\term{r} : \typeformer{itree}~\typeformer{E}~\type{R}}{r : \type{R}}
\end{equation}
\begin{equation}
  \infer[\mathtt{I}_{\function{Vis}}]{\function{Vis}~\term{a}~\term{f} : \typeformer{itree}~\typeformer{E}~\type{R}}{\type{A} : \universe{\mathcal{U}} & \term{a} : \typeformer{E}~\type{A} & \term{f} : \type{A} \rightarrow \typeformer{itree}~\typeformer{E}~\type{R}}.
\end{equation}
Elimination rules
\begin{equation}
  \infer[\mathtt{E}_{\function{Tau}}]{\function{Tau}~\term{t} : \typeformer{itree}~\typeformer{E}~\type{R}}{\term{t} : \typeformer{itree}~\typeformer{E}~\type{R}}.
\end{equation}

\subsection{Delay Monad}
We start by looking at itrees without the \(\function{Vis}\) constructor, this type is also know as the delay monad\todo{check this statement}. We construct this type by letting \(\container{S} = (\unit + \type{R},\lambda \{ \texttt{inl}~\_ \rightarrow \unit~;~\texttt{inr}~\_ \rightarrow \empt \})\), we then get the polynomial functor
\begin{equation}
  \functor{P}_{\container{S}}(\type{X}) = \sum_{\term{x} : \unit + \type{R}} \lambda \{ \texttt{inl}~\_ \rightarrow \unit ; \texttt{inr}~\_ \rightarrow \empt\}~\term{x} \rightarrow \type{X},
\end{equation}
which is equal to
\begin{equation}
  \functor{P}_{\container{S}}(\type{X}) = \type{X} + \type{R} \times (\empt \rightarrow \type{X}).
\end{equation}
We know that \(\empt \rightarrow \type{X} \equiv \unit\), so we can reduce further to
\begin{equation}
  \functor{P}_{\container{S}}(\type{X}) = \type{X} + \type{R}
\end{equation}
meaning we get the following diagram.

\begin{figure}[h]
  \centering
  \begin{tikzcd}[row sep=large,column sep=large]
    \type{R} \ar[dr,"\function{now}"'] \ar[r,"\texttt{inl}"] & \type{R} + \typeformer{M}  & \typeformer{M} \ar[l,"\texttt{inr}"'] \ar[dl,"\function{later}"] \\
    & \typeformer{M} \ar[u,"\function{out}"']
  \end{tikzcd}
  \caption{Delay monad}
\end{figure}
\noindent Meaning we can define the operations \(\function{now}\) and \(\function{later}\) using \(\function{in} = \function{out}^{-1}\) together with the injections \(\mathtt{inl}\) and \(\mathtt{inr}\). \todo{(Later = Tau, Ret = Now)}

\subsection{Tree}
Now lets look at the example where we remove the \(\function{Tau}\) constructor. We let
\begin{equation}
\container{S} = \left(\type{R} + \sum_{\type{A} : \universe{\mathcal{U}}} \typeformer{E}~\type{A}~,~\lambda \{ \texttt{inl}~\_ \rightarrow \empt ~;~ \texttt{inr}~(\type{A},\term{e}) \rightarrow \type{A} \} \right).
\end{equation}
This will give us the polynomial functor
\begin{equation}
  \functor{P}_{\container{S}}(\type{X}) = \sum_{\term{x} : \type{R} + \sum_{\type{A} : \universe{\mathcal{U}}} \typeformer{E}~\type{A}} \lambda \{ \texttt{inl}~\_ \rightarrow \empt ~;~ \texttt{inr}~(\type{A},\term{e}) \rightarrow \type{A} \}~\term{x} \rightarrow \type{X},
\end{equation}
which simplifies to
\begin{equation}
  \functor{P}_{\container{S}}(\type{X}) = (\type{R} \times (\empt \rightarrow \type{X})) + (\sum_{\type{A} : \universe{\mathcal{U}}} \typeformer{E}~\type{A} \times (\type{A} \rightarrow \type{X})),
\end{equation}
and further
\begin{equation}
  \functor{P}_{\container{S}}(\type{X}) = \type{R} + \sum_{\type{A} : \universe{\mathcal{U}}} \typeformer{E}~\type{A} \times (\type{A} \rightarrow \type{X}).
\end{equation}
We get the following diagram for the \(\functor{P}\)-coalgebra.

\begin{figure}[h]
  \centering
  \begin{tikzcd}[row sep=large,column sep=large]
    \type{R} \ar[dr,"\function{Ret}"'] \ar[r,"\texttt{inl}"] & \type{R} + \sum_{\type{A} : \universe{U}} \typeformer{E}~\type{A} \times (\type{A} \rightarrow \typeformer{M})  & \sum_{\type{A} : \universe{U}} \typeformer{E}~\type{A} \times (\type{A} \rightarrow \typeformer{M}) \ar[l,"\texttt{inr}"'] \ar[dl,"\function{Vis}"] \\
    & \typeformer{M} \ar[u,"\function{out}"]
  \end{tikzcd}
  \caption{TODO}
\end{figure}
\noindent Again we can define \(\function{Ret}\) and \(\function{Vis}\) using the \(\function{in}\) function.

\subsection{ITrees}
Now we should have all the knowledge needed to make ITrees using \texttt{M}-types. We define ITrees by the container
\begin{equation}
  \container{S} = \left( \unit+\type{R}+\sum_{\type{A}:\universe{\mathcal{U}}} (\typeformer{E}~\type{A}) ~~,~~ \lambda\left\{ \texttt{inl}~(\texttt{inl}~\_) \rightarrow \unit ~;~ \texttt{inl}~(\texttt{inr}~\_) \rightarrow \empt ~;~ \texttt{inr} (\type{A},\_)  \rightarrow \type{A} \right\} \right).
\end{equation}
Such that the (reduced) polynomial functor becomes
\begin{equation}
  \functor{P}_{\container{S}}(\type{X}) = \type{X} + \type{R} + \sum_{\type{A} : \universe{\mathcal{U}}} ((\typeformer{E}~\type{A}) \times (\type{A} \rightarrow \type{X}))
\end{equation}
Giving us the diagram
\begin{figure}[h]
  \centering
  \begin{tikzcd}[row sep=large,column sep=large]
    \type{R} \ar[ddrr,"\function{Ret}"'] \ar[r,"\texttt{inl}\,\circ\,\texttt{inr}"] & \typeformer{M} + \type{R} + \sum_{\type{A} : \universe{\mathcal{U}}} (\typeformer{E}~\type{A} \times (\type{A} \rightarrow \typeformer{M})) & & \sum_{\type{A} : \universe{U}} \typeformer{E}~\type{A} \times (\type{A} \rightarrow \typeformer{M}) \ar[ll,"\texttt{inr}"'] \ar[ddl,"\function{Vis}"] \\
    & & \typeformer{M} \ar[ul,"\texttt{inl}\,\circ\,\texttt{inl}"'] \ar[d,"\function{Tau}"'] \\
    & & \typeformer{M} \ar[uul,"\function{out}"]
  \end{tikzcd}
  \caption{TODO}
\end{figure}

\section{Co-induction Principle for M-types}
We can now construct a bisimulation: forall coalgebras \(\coalg{\type{C}}{\term{\gamma}} : \typeformer{Coalg}_{\container{S}}\), if we have a relation \(\typeformer{\mathcal{R}} : \type{C} \rightarrow \type{C} \rightarrow \universe{\mathcal{U}}\), and a type \(\type{\overline{\mathcal{R}}} = \sum_{\term{a} : \type{C}} \sum_{\term{b} : \type{C}} \typeformer{\mathcal{R}}~\term{a}~\term{b}\), such that \(\type{\overline{\mathcal{R}}}\) and \(\term{\alpha}_{\type{\overline{\mathcal{R}}}} : \type{\overline{\mathcal{R}}} \rightarrow \functor{P}(\type{\overline{\mathcal{R}}})\) forms a \(\functor{P}\)-coalgebra \(\coalg{\type{\overline{\mathcal{R}}}}{\term{\alpha}_{\type{\overline{\mathcal{R}}}}} : \mathtt{Coalg}_{\container{S}}\), making the diagram in \figref{fig:coalgebra-coinduction-definition} commute (where \(\Longrightarrow\) are \(\functor{P}\)-coalgebra morphisms).
\begin{figure}[h]
  \centering
  \begin{tikzcd}[row sep=huge,column sep=huge]
    \coalg{\type{C}}{\term{\gamma}} & \coalg{\type{\overline{\mathcal{R}}}}{\term{\alpha}_{\type{\overline{\mathcal{R}}}}} \ar[l,Rightarrow,"{\pi_1}^{\type{\overline{\mathcal{R}}}}"'] \ar[r,Rightarrow,"{\pi_2}^{\type{\overline{\mathcal{R}}}}"] & \coalg{\type{C}}{\term{\gamma}}
  \end{tikzcd}
  \caption{Bisimulation for a coalgebra}
  \label{fig:coalgebra-coinduction-definition}
\end{figure}
\noindent Furthermore for any bisimulation over a final \(\functor{P}\)-coalgebra \(\coalg{\typeformer{M}}{\term{\mathtt{out}}}:\mathtt{Coalg}_{\container{S}}\) we get the diagram in \figref{fig:final-coalgebra-coinduction},
\begin{figure}[h]
  \centering
  \begin{tikzcd}[row sep=huge,column sep=huge]
    \coalg{\typeformer{M}}{\function{out}} & \coalg{\type{\overline{\mathcal{R}}}}{\term{\alpha}_{\type{\overline{\mathcal{R}}}}} \ar[l,Rightarrow,"{\pi_1}^{\type{\overline{\mathcal{R}}}}"'] \ar[r,Rightarrow,"{\pi_2}^{\type{\overline{\mathcal{R}}}}"] & \coalg{\typeformer{M}}{\function{out}}
  \end{tikzcd}
  \caption{Bisimulation principle for final coalgebra}
  \label{fig:final-coalgebra-coinduction}
\end{figure}
\noindent where \(\pi_1^{\type{\overline{\mathcal{R}}}} = ~! = \pi_2^{\type{\overline{\mathcal{R}}}}\), which means given \(\term{r} : \term{\mathcal{R}}(\term{m},\term{m'})\) we get \(\term{m} = \pi_1^{\type{\overline{\mathcal{R}}}} (\term{m}, \term{m'}, \term{r}) = \pi_2^{\type{\overline{\mathcal{R}}}} (\term{m}, \term{m'}, \term{r}) = \term{m'}\).
\\ \\
\todo[inline]{Better introduction to this proof!}
We want to define a co-induction principle from any bisimulation relation over a final coalgebra, that is if \(\term{R}\) gives a bisimulation, then it is true that
\begin{equation}
  \term{R} \equiv \term{\equiv}
\end{equation}
meaning we can use the relation \(\term{R}\), to show that two things of an \texttt{M}-type are equivalent. So we want to construct an isomorphism between \(\term{R}\) and the equivalence relation \(\equiv\), to do this we must construct functions
\begin{align}
  \term{p} &: \term{R} \rightarrow \term{\equiv} \\
  \term{q} &: \term{\equiv} \rightarrow \term{R}
\end{align}
and relations
\begin{align}
  \term{\alpha} &: \term{p} \circ \term{q} \equiv \term{\mathtt{id}_{\equiv}} \\
  \term{\beta} &: \term{q} \circ \term{p} \equiv \term{\mathtt{id}_{R}}
\end{align}
\todo[inline]{Complete the construction of equality from any bisimulation relation over an M-type}

\subsection{Bisimulation of Streams}
\todo[inline]{TODO}
\subsection{Bisimulation of Delay monad}
\todo[inline]{TODO}
\subsection{Bisimulation of ITrees}
We define our bisimulation coalgebra from the strong bisimulation relation \(\mathcal{R}\), defined by the following rules.
\begin{equation}
  \infer[\mathtt{EqRet}]{\mathtt{Ret}~\term{a} \cong \mathtt{Ret}~\term{b}}{\term{a},\term{b} : \type{R} & \term{a} \equiv_{\type{R}} \term{b}}
\end{equation}
\begin{equation}
  \infer[\mathtt{EqTau}]{\mathtt{Tau}~\term{t} \cong \mathtt{Tau}~\term{u}}{\term{t},\term{u} : \mathtt{itree}~\typeformer{E}~\type{R} & \term{t} \cong \term{u}}
\end{equation}
\begin{equation}
  \infer[\mathtt{EqVis}]{\mathtt{Vis}~\term{e}~\term{k_1} \cong \mathtt{Tau}~\term{e}~\term{k_2}}{\type{A} : \universe{\mathcal{U}} & \term{e} : \typeformer{E}~\type{A} & \term{k_1},\term{k_2} : \type{A} \rightarrow \mathtt{itree}~\typeformer{E}~\type{R} & \term{t} \cong \term{u}}
\end{equation}
Now we just need to define \(\term{\alpha}_{\type{\overline{\mathcal{R}}}}\) \todo{define the \(\alpha_{\overline{R}}\) function}. Now we have a bisimulation relation, which is equivalent to equality, using what we showed in the previous section.

\section{Stream Formalization using M-types}
When taking the head of a stream, we get
\begin{align}
  \function{hd}~(\function{cons}~x~xs) &\equiv \pi_1~\function{out}~(\function{cons}~x~xs) \\
                                       &\equiv \pi_1~\function{out}~(\function{in}~(x , \lambda\,\_,xs)) \\
                                       &\equiv \pi_1~(x , \lambda\,\_,xs) \\
                                       &\equiv x
\end{align}


\section{Quotient M-type}
Since we know that M-types preserves the H-level, we can use set-truncated quotients, to define quotient M-types, for examples we can define weak bisimulation of the delay monad as
\todo[inline]{Quotients of the delay monad}

\section{Closure properties of M-types}
We define the product of two containers
\begin{equation}
  \containerpair{\type{A}}{\type{B}} \times \containerpair{\type{C}}{\type{D}} \equiv \containerpair{\type{A} \times \type{C}}{\lambda\,(\term{a} , \term{c}), \typeformer{B}~\term{a} \times \typeformer{D}~\term{c}},
\end{equation}
we can lift this rule, through the diagram in \figref{fig:container-products}, used to define M-types.
\begin{figure}[h]
  \centering
  \begin{tikzcd}[row sep=huge,column sep=huge]
    \type{A} & \type{A} \times \type{C} \ar[l,"\pi_1"'] \ar[r,"\pi_2"] & \type{C} \\
    \functor{P} (\typeformer{M}_{\containerpair{\type{A}}{\type{B}}}) \ar[u,"\pi_1"'] & \functor{P}(\typeformer{M}_{\containerpair{\type{A}}{\type{B}} \times \containerpair{\type{C}}{\type{D}}}) \ar[u,"\pi_1"'] & \functor{P} (\typeformer{M}_{\containerpair{\type{C}}{\type{D}}}) \ar[u,"\pi_1"'] \\
    \typeformer{M}_{\containerpair{\type{A}}{\type{B}}} \ar[u,"\function{out}"'] \ar[d,"\pi_{(!)}"] & \typeformer{M}_{\containerpair{\type{A}}{\type{B}} \times \containerpair{\type{C}}{\type{D}}} \ar[u,"\function{out}"'] \ar[d,"\pi_{(!)}"] & \typeformer{M}_{\containerpair{\type{C}}{\type{D}}} \ar[u,"\function{out}"'] \ar[d,"\pi_{(!)}"] \\
    {\functor{P}_{\containerpair{\type{A}}{\type{B}}}}^n~\unit \ar[d,"\pi_{(n)}"] & {\functor{P}_{\containerpair{\type{A}}{\type{B}} \times \containerpair{\type{C}}{\type{D}}}}^n~\unit \ar[l] \ar[r] \ar[d,"\pi_{(n)}"] & {\functor{P}_{\containerpair{\type{C}}{\type{D}}}}^n~\unit \ar[d,"\pi_{(n)}"] \\
    \vdots \ar[d,"\pi_{(1)}"] & \vdots \ar[d,"\pi_{(1)}"] & \vdots \ar[d,"\pi_{(1)}"] \\
    \functor{P}_{\containerpair{\type{A}}{\type{B}}}~\unit \ar[d,"\pi_{(0)}"] & \functor{P}_{\containerpair{\type{A}}{\type{B}} \times \containerpair{\type{C}}{\type{D}}}~\unit \ar[l] \ar[r] \ar[d,"\pi_{(0)}"] & \functor{P}_{\containerpair{\type{C}}{\type{D}}}~\unit \ar[d,"\pi_{(0)}"] \\
    \unit & \unit \ar[l,"\texttt{refl}"] \ar[r,"\texttt{refl}"'] & \unit
  \end{tikzcd}
  \caption{TODO}
  \label{fig:container-products}
\end{figure}
We now prove that
\begin{equation}
  {\functor{P}_{\containerpair{\type{A}}{\type{B}}}}^n~\unit \times {\functor{P}_{\containerpair{\type{C}}{\type{D}}}}^n~\unit \equiv {\functor{P}_{\containerpair{\type{A}}{\type{B}} \times \containerpair{\type{C}}{\type{D}}}}^n~\unit,
\end{equation}
by induction on \(n\). For \(n = 0\), we have \(\unit \times \unit \equiv \unit\). For \(n = m + 1\), we may assume
\begin{equation}
{\functor{P}_{\containerpair{\type{A}}{\type{B}}}}^m~\unit \times {\functor{P}_{\containerpair{\type{C}}{\type{D}}}}^m~\unit \equiv {\functor{P}_{\containerpair{\type{A}}{\type{B}} \times \containerpair{\type{C}}{\type{D}}}}^m~\unit,
\end{equation}
and show 
\begin{align}
  &~ {\functor{P}_{\containerpair{\type{A}}{\type{B}}}}^{m+1}~\unit \times {\functor{P}_{\containerpair{\type{C}}{\type{D}}}}^{m+1}~\unit \\
   \equiv&~ \functor{P}_{\containerpair{\type{A}}{\type{B}}} ({\functor{P}_{\containerpair{\type{A}}{\type{B}}}}^m~\unit) \times \functor{P}_{\containerpair{\type{C}}{\type{D}}} ({\functor{P}_{\containerpair{\type{C}}{\type{D}}}}^m~\unit) \\
   \equiv&~ \sum_{\term{a} : \type{A}} \type{B}~\term{a} \rightarrow {\functor{P}_{\containerpair{\type{A}}{\type{B}}}}^m~\unit \times \sum_{\term{c} : \type{C}} \type{D}~\term{c} \rightarrow {\functor{P}_{\containerpair{\type{C}}{\type{D}}}}^m~\unit \\
  \equiv&~ \sum_{\term{a} , \term{c} : \type{A} \times \type{C}} (\type{B}~\term{a} \rightarrow {\functor{P}_{\containerpair{\type{A}}{\type{B}}}}^m~\unit) \times (\type{D}~\term{c} \rightarrow {\functor{P}_{\containerpair{\type{C}}{\type{D}}}}^m~\unit) \\
  \equiv&~ \sum_{\term{a} , \term{c} : \type{A} \times \type{C}} \type{B}~\term{a} \times \type{D}~\term{c} \rightarrow {\functor{P}_{\containerpair{\type{A}}{\type{B}}}}^m~\unit \times {\functor{P}_{\containerpair{\type{C}}{\type{D}}}}^m~\unit \\
  \equiv&~ \sum_{\term{a} , \term{c} : \type{A} \times \type{C}} \type{B}~\term{a} \times \type{D}~\term{c} \rightarrow {\functor{P}_{\containerpair{\type{A}}{\type{B}} \times \containerpair{\type{C}}{\type{D}}}}^m~\unit \\
  \equiv&~ \functor{P}_{\containerpair{\type{A}}{\type{B}} \times \containerpair{\type{C}}{\type{D}}} ({\functor{P}_{\containerpair{\type{A}}{\type{B}} \times \containerpair{\type{C}}{\type{D}}}}^m~\unit) \\
  \equiv&~ {\functor{P}_{\containerpair{\type{A}}{\type{B}} \times \containerpair{\type{C}}{\type{D}}}}^{m+1}~\unit
\end{align}
taking the limit we get
\begin{equation}
  \typeformer{M}_{\containerpair{\type{A}}{\type{B}}} \times \typeformer{M}_{\containerpair{\type{C}}{\type{D}}} \equiv \typeformer{M}_{\containerpair{\type{A}}{\type{B}} \times \containerpair{\type{C}}{\type{D}}}
\end{equation}
as an example hereof lets look at streams, where we actually get
\begin{equation}
  \typeformer{stream}~\type{A} \times \typeformer{stream}~\type{B} \equiv \typeformer{M}_{\containerpair{\type{A}}{\lambda~\_,\unit}} \times \typeformer{M}_{\containerpair{\type{B}}{\lambda~\_,\unit}} \equiv \typeformer{M}_{\containerpair{\type{A}}{\lambda~\_,\unit} \times \containerpair{\type{B}}{\lambda~\_,\unit}} \equiv \typeformer{stream}~(\type{A} \times \type{B}) \label{eq:zip-equality}
\end{equation}
as expected, transporting along equality \eqref{eq:zip-equality} gives us the definition for \function{zip}.

\subsection{Zip Function}
We want the diagram in \figref{fig:zip-computation-rules} to commute,
\begin{figure}[h]
  \centering
  \begin{tikzcd}[row sep=huge,column sep=huge]
    & A \times B & \\
    \typeformer{stream}~\type{A} \times \typeformer{stream}~\type{B} \ar[ur,"\function{hd} \times \function{hd}"'] \ar[rr,"\function{zip}"] \ar[d,"\function{tl} \times \function{tl}"] & & \typeformer{stream}~(\type{A} \times \type{B}) \ar[ul,"\function{hd}"] \ar[d,"\function{tl}"] \\
    \typeformer{stream}~\type{A} \times \typeformer{stream}~\type{B} \ar[rr,"\function{zip}"] & & \typeformer{stream}~(\type{A} \times \type{B})
  \end{tikzcd}
  \caption{TODO}
  \label{fig:zip-computation-rules}
\end{figure}
meaning we get the computation rules
\begin{equation}
  (\function{hd} \times \function{hd}) \equiv \function{hd} \circ \function{zip}
\end{equation}
\begin{equation}
  \function{zip} \circ (\function{tl} \times \function{tl}) \equiv \function{tl} \circ \function{zip}
\end{equation}
we can define the zip function as we did in the end of the last section. Another way to define the zip function is more directly, using the following lifting property of M-types
\begin{equation}
\begin{aligned}
  &\function{lift_M}~\left(x : \prod_{n : \mathbb{N}} (\type{A} \rightarrow {\functor{P}_{\container{S}}}^n~\unit)\right)~\left(u : \prod_{n : \mathbb{N}} (\type{A} \rightarrow \pi_n (x_{n+1} a) \equiv x_n a)\right) (a : \type{A}) : \function{M}~\container{S} := \\
  &\qquad\qquad (\lambda\,n, x~n~a) , (\lambda\,n~i, p~n~a~i).
\end{aligned} \label{eq:M-lift}
\end{equation}
To use this definition, we first define some helper functions
\begin{equation}
  \function{zip}_{X}~n~(x , y) = \begin{cases} \unit & \mathtt{if}~n = 0 \\ (\function{hd}~x , \function{hd}~y) , (\lambda\,\_, \function{zip}_{X}~m~(\function{tl}~x,\function{tl}~y)) , & \mathtt{if}~n = m + 1 \end{cases}
\end{equation}
\begin{equation}
  \function{zip}_{\pi}~n~(x , y) = \begin{cases} \function{refl} & \mathtt{if}~n = 0 \\ \lambda\,i,(\function{hd}~x , \function{hd}~y) , (\lambda\,\_, \function{zip}_{\pi}~m~(\function{tl}~x,\function{tl}~y)~i) , & \mathtt{if}~n = m + 1 \end{cases},
\end{equation}
we can then define
\begin{equation}
  \function{zip}_{lift}~(x , y) := \function{lift_M}~\function{zip}_X~\function{zip}~(x,y).
\end{equation}
\subsubsection{Equality of Zip Definitions}
We would expect that the two definitions for zip are equal
\begin{align}
  \texttt{transport}_{?}~a &\equiv \function{zip}_{lift}~a \\
                           &\equiv \function{lift_M}~\function{zip}_X~\function{zip}_\pi~(x,y) \\
                           &\equiv (\lambda\,n, \function{zip}_X~n~(x,y)) , (\lambda\,n~i, \function{zip}_\pi~n~(x,y)~i)
\end{align}
zero case \(X\)
\begin{equation}
  \function{zip}_X~0~(x,y) \equiv \unit
\end{equation}
Sucessor case \(X\)
\begin{align} 
  \function{zip}_X~(m+1)~(x,y) &\equiv (\function{hd}~x, \function{hd}~y),(\lambda\,\_,\function{zip}_X~m~(\function{tl}~x,\function{tl}~y)) \\
   &\equiv (\function{hd}~x, \function{hd}~y),(\lambda\,\_, {?}~(\function{tl}~a)) \\
   &\equiv (\function{hd}~(\texttt{transport}_{?} a)),(\lambda\,\_, \texttt{transport}_{?} (\function{tl}~a)) \\
   &\equiv transport_? a \\
\end{align}
Zero case \(\pi\): \((\lambda\,i, \function{zip}_\pi~0~(x,y)~i \equiv refl)\).
\begin{align}
  &\equiv () , (\lambda\,i, \function{zip}_\pi~0~(x,y)~i) \\
  &\equiv \unit , \function{refl} \\
\end{align}
successor case
\begin{align}
  &\equiv (\function{zip}_X~(m+1)~(x,y)) , (\lambda\,i, \function{zip}_\pi~(m+1)~(x,y)~i) \\
  &\equiv ((\function{hd}~x, \function{hd}~y),(\lambda\,\_,\function{zip}_X~m~(\function{tl}~x,\function{tl}~y))) , (\lambda\,i, (\function{hd}~x, \function{hd}~y),(\lambda\,\_,\function{zip}_\pi~m~(\function{tl}~x,\function{tl}~y)~i))
\end{align}
\todo[inline]{Complete this proof}

\subsection{Examples of Fixed Points}
\subsubsection{Zeros}
Let us try to define the zero stream, we do this by lifting the functions
\begin{equation}
  \function{const_X}~(n : \type{\mathbb{N}})~(c : \type{\mathbb{N}})  := \begin{cases} \unit & n = 0 \\ (c , \lambda\,\_, \function{const_X}~m~c) & n = m + 1 \end{cases}
\end{equation}
\begin{equation}
  \function{const_\pi}~(n : \type{\mathbb{N}})~(c : \type{\mathbb{N}}) := \begin{cases} \mathtt{refl} & n = 0 \\ \lambda\,i,(c , \lambda\,\_, \function{const_\pi}~m~c~i) & n = m + 1 \end{cases}
\end{equation}
to get the definition of zero stream
\begin{equation}
  \mathtt{zeros} := \function{lift_M}~\function{const_X}~\function{const_\pi}~0.
\end{equation}
We want to show that we get the expected properties, such as
\begin{align}
  \function{hd}~\mathtt{zeros} &\equiv 0 \\
  \function{tl}~\mathtt{zeros} &\equiv \mathtt{zeros}
\end{align}

\subsubsection{Spin}
We want to define spin, as being the fixed point \(\mathtt{spin} = \mathtt{later}~\mathtt{spin}\), so that is again a final coalgebra, but of a M-type (which is a final coalgebra)

\begin{figure}[h]
  \centering
  \begin{tikzcd}[row sep=huge,column sep=huge]
    \function{spin} \ar[d,"\function{later}"] \\
    \function{later}~\function{spin}
  \end{tikzcd}
  \caption{TODO}
\end{figure}
\noindent Since it is final, it also must be unique, meaning that there is just one program that spins forever, without returning a value, meaning every other program must return a value. If we just


%%%%%%%%%%%%%%%%%%%%%%%%%%%%%%%%%%%%%%%%%%%%%%%%%%%%%%%%%%%%%%%%%%%%%%%

\chapter{Additions to the Cubical Agda Library}
\section{Lemma 10}
\texttt{M}-types are part of a final coalgebra, formally \(\forall\,\container{S}~\coalg{\type{C}}{\term{\gamma}}, (\coalg{\type{C}}{\term{\gamma}} \Rightarrow \coalg{\typeformer{M}}{\term{\mathtt{out}}}) \equiv \unit\)
\begin{align}
  U &\equiv \sum_{f : C \rightarrow \mathcal{L}} \mathtt{out} \circ f \equiv \mathtt{step}~f \\
    &\equiv \sum_{f : C \rightarrow \mathcal{L}} \mathtt{in} \circ \mathtt{out} \circ f \equiv \mathtt{in} \circ \mathtt{step}~f\\
    &\equiv \sum_{f : C \rightarrow \mathcal{L}} f \equiv \mathtt{in} \circ \mathtt{step}~f\\
    &\equiv \sum_{f : C \rightarrow \mathcal{L}} f \equiv \Psi~f\\
    &\equiv \sum_{c : Cone} e~c \equiv \Psi~(e~c)\\
    &\equiv \sum_{c : Cone} e~c \equiv e~(\phi~c) \\
    &\equiv \sum_{c : Cone} c \equiv \phi~c \\
    &\equiv \sum_{(u,q) : Cone} (u , q) \equiv (\phi_0~u , \phi_1~u~q) \\
    &\equiv \sum_{(u,q) : Cone} \sum_{p : u \equiv \phi_0~u} q \equiv_{\lambda i, Cone_1 (p~i)} \phi_1~u~q \\
    &\equiv \sum_{(u,p) : \sum_{u : Cone_0} u \equiv \phi_0~u} \sum_{q : Cone_1~u} q \equiv_{\lambda i, Cone_1 (p~i)} \phi_1~u~q \\
    &\equiv \sum_{q : Cone_1 u_0} q \equiv_{\lambda i, Cone_1 (funExt~p_0~i)} \phi_1~u_0~q \\
    &\vdots \\
    &\equiv \unit
\end{align}

\section{Lemma 13}
\begin{align}
  &~\sum_{((a,q):\sum_{a : \mathbb{N} \rightarrow A} \prod_{n : \mathbb{N}} a_{n+1} \equiv a_n)} \sum_{(u : \prod_{n : \mathbb{N}} B a_n \rightarrow X_n)} \prod_{(n : \mathbb{N})} \pi \circ u_{n + 1} \equiv_{\lambda i, B (q_n~i) \rightarrow X_n} u_n \\
  \equiv&~\sum_{a : A} \sum_{u : \prod_{n : \mathbb{N}} B a \rightarrow X_n} \prod_{n : \mathbb{N}} \pi \circ u_{n+1} \equiv u_n \\
  \equiv&~\sum_{a : A} \sum_{u : \prod_{n : \mathbb{N}} B a \rightarrow X_n} \prod_{n : \mathbb{N}} \pi \circ u_{n+1} \equiv u_n
\end{align}

\section{Missing postulates}
\subsubsection{Combine}
For all \(\term{X} : \type{\mathbb{N}} \rightarrow \universe{U}\) and \(p : \prod_{n : \type{\mathbb{N}}} \term{X}~(n + 1) \rightarrow \term{X}~n \rightarrow \universe{U}\)
\begin{align}
  &~\sum_{\term{x_0} : \type{X_0}} \sum_{\term{y} : \prod_{\term{n} : \type{\mathbb{N}}} \type{X_{n+1}}} (\term{p}~\type{y_0}~\type{x_0}) \times \left(\prod_{\term{n} : \type{\mathbb{N}}} \term{p}~\type{y_{n + 1}}~\type{y_n}\right) \\
  \equiv&~\sum_{\term{x_0} : \type{X_0}} \sum_{\term{y} : \prod_{\term{n} : \type{\mathbb{N}}} \type{X_{n+1}}} (\term{p}~\type{y_0}~\type{x_0}) \times \left(\prod_{\term{n} : \type{\mathbb{N}}} \term{p}~\type{y_{n + 1}}~\type{y_n}\right) \\
  \equiv&~\sum_{\term{x} : \prod_{\term{n} : \type{\mathbb{N}}} \rightarrow \type{X_n}} (\term{p}~\type{x_1}~\type{x_0}) \times \left(\prod_{\term{n} : \type{\mathbb{N}}} → \term{p}~\type{x_{n+2}}~\type{x_{n+1}}\right)
\end{align}

\chapter{Conclusion}
\label{ch:conclusion}

\todo[inline]{conclude on the problem statement from the introduction}

%%%%%%%%%%%%%%%%%%%%%%%%%%%%%%%%%%%%%%%%%%%%%%%%%%%%%%%%%%%%%%%%%%%%%%%

\cleardoublepage
\addcontentsline{toc}{chapter}{Bibliography}
\bibliographystyle{plain} 
\bibliography{Thesis}

%%%%%%%%%%%%%%%%%%%%%%%%%%%%%%%%%%%%%%%%%%%%%%%%%%%%%%%%%%%%%%%%%%%%%%%

\cleardoublepage
\appendix
\chapter{The Technical Details}

\todo[inline]{\dots}

\end{document}

